% Part: incompleteness
% Chapter: arithmetization-syntax
% Section: coding-formulas

\documentclass[../../../include/open-logic-section]{subfiles}

\begin{document}

\olfileid[hi]{inc}{art}{frm}
\olsection{\printtoken{P}{formula} का कूटलेखन}

संबंध $\fn{Term}(x)$ को आदिम पुनरावर्ती रूप से परिभाषित कर लेने के बाद हम
उसका उपयोग करके !!{formula} के लिए संगत संबंध $\fn{Frm}(x)$ को आदिम
पुनरावर्ती रूप से परिभाषित कर सकते हैं।

\begin{prop}
संबंध $\fn{Atom}(x)$, जो तभी लागू होता है जब $x$ किसी परमाणु !!{formula} की
ग्योडेल (G\"odel) संख्या हो, आदिम पुनरावर्ती है।
\end{prop}

\begin{proof}
संख्या $x$ किसी परमाणु !!{formula} की ग्योडेल (G\"odel) संख्या है यदि और
केवल यदि निम्न में से कोई एक बात लागू होती है:
\begin{enumerate}
\item ऐसे $n$, $j < x$ और $z < x$ हैं कि प्रत्येक $i < n$ के लिए
  $\fn{Term}((z)_i)$ और $x = $
\[
\Gn{\Obj P^n_j(} \concat \fn{flatten}(z) \concat \Gn{)}.
\]
\item ऐसे $z_1, z_2 < x$ हैं कि $\fn{Term}(z_1)$, $\fn{Term}(z_2)$ और $x = $
\[
\Gn{{\eq}(} \concat z_1 \concat \Gn{,} \concat z_2 \concat{\Gn{)}}.
\]
  \tagitem{prvFalse}{$x = \Gn{\lfalse}$.}{}
  \tagitem{prvTrue}{$x = \Gn{\ltrue}$.}{}
\end{enumerate}
\end{proof}

\begin{prop}
\ollabel{prop:frm-primrec}
संबंध $\fn{Frm}(x)$, जो तभी लागू होता है जब $x$ किसी !!a{formula} की ग्योडेल (G\"odel)
 संख्या हो, आदिम पुनरावर्ती है।
\end{prop}

\begin{proof}
चिह्नों का अनुक्रम~$s$ !!a{formula} है यदि और केवल यदि !!{formula} का कोई
निर्माण-अनुक्रम~$s_0$, \dots, $s_{k-1} = s$ यह दर्ज करता है कि निर्माण-नियमों
के अनुसार परमाणु !!{formula} से~$s$ कैसे बना। प्रत्येक $s_i$ का कूट (और
वास्तव में अनुक्रम $\tuple{s_0, \dots, s_{k-1}}$ का कूट भी),~$s$ के कूट~$x$
से छोटा है।
\end{proof}

\begin{prob}
\olref[inc][art][trm]{prop:term-primrec} के पहले प्रमाण की पद्धति पर
\olref[inc][art][frm]{prop:frm-primrec} का विस्तृत प्रमाण दें।
\end{prob}

\begin{prop}
\ollabel{prop:freeocc-primrec}
संबंध $\fn{FreeOcc}(x, z, i)$, जो तभी लागू होता है जब ग्योडेल (G\"odel)
संख्या~$x$ वाले सूत्र का $i$-वाँ चिह्न ग्योडेल (G\"odel) संख्या~$z$ वाले चर
की मुक्त उपस्थिति हो, आदिम पुनरावर्ती है।
\end{prop}

\begin{proof}
अभ्यास।
\end{proof}

\begin{prob}
\olref[inc][art][frm]{prop:freeocc-primrec} सिद्ध करें। आप इस तथ्य का उपयोग
कर सकते हैं कि किसी !!a{formula} की प्रत्येक ऐसी उपडोरी जो स्वयं !!a{formula}
है, उसका उप-!!{formula} है।
\end{prob}

\begin{prop}
गुण $\fn{Sent}(x)$, जो तभी लागू होता है जब $x$ किसी !!a{sentence} की ग्योडेल (G\"odel)
 संख्या हो, आदिम पुनरावर्ती है।
\end{prop}

\begin{proof}
!!{sentence} ऐसा !!a{formula} है जिसमें !!{variable} की कोई मुक्त उपस्थिति
नहीं होती। अतः $\fn{Sent}(x)$ लागू होता है यदि और केवल यदि
\begin{multline*}
\bforall{i<\len{x}}{\bforall{z<x}{}}\\
(\bexists{j<z}{z=\Gn{\Obj v_j}} \lif \lnot\fn{FreeOcc}(x,z,i)).
\end{multline*}
\end{proof}


\end{document}
